% =============================================================================
% Step-6 (unweighted origin) supplement: 2 archs x 5 losses x 3 solvers x 3 groups,
% multi-molecule training with V_xc-aware losses.
% Compile alongside main.tex (shared references.bib).
% =============================================================================
\documentclass[11pt,letterpaper]{article}
\usepackage[margin=1in]{geometry}
\usepackage{graphicx, booktabs, amsmath, amssymb}
\usepackage[hidelinks]{hyperref}
\usepackage{natbib, xcolor, microtype, caption}
\captionsetup{font=small, labelfont=bf}
\graphicspath{
  {../step6_unweighted_analysis/figures/}
  {../../notebooks/checkpoints_step6/unweighted/figures/}
}

\title{Step-6 Supplement (unweighted origin): Multi-Molecule Training with
  $V_{xc}$-Aware Losses}
\author{Alec P.\ Wills}
\date{\today}

\begin{document}
\maketitle

\section*{Run scope}
90 specs = 2 architectures (\texttt{deep\_combined}, \texttt{deep\_combined\_attn})
$\times$ 5 losses (\texttt{L1\_B}, \texttt{L2\_C\_anchor}, \texttt{L3\_balanced\_vxc},
\texttt{L4\_balanced\_vxc\_anchor}, \texttt{L5\_gradnorm\_vxc})
$\times$ 3 solvers (\texttt{oneshot}, \texttt{fixed\_j\_3}, \texttt{full\_3})
$\times$ 3 training groups (group1: H$_2$O short; group2: H$_2$O+C$_2$H$_2$ short;
group3: H$_2$O+C$_2$H$_2$ long). All specs use the unweighted-pretrain origin
($w_x \!=\! w_c \!=\! 1$ inside \texttt{xcquinox/alec/pretrain.py}).

\section{Headline numbers}

\begin{table}[h]
\centering
\caption{Step-6 unweighted-origin headlines (verified against
\texttt{headline\_stats.json}).}
\small
\begin{tabular}{lr}
\toprule
metric & value \\
\midrule
\# specs & 90 \\
best AE-MAE (kcal/mol) & 0.00793 \\
best AE spec & \texttt{group3/deep\_combined\_attn/L2\_C\_anchor/oneshot} \\
best $\rho$-RMSE (e/bohr$^3$) & $1.32\!\times\!10^{-4}$ \\
best $\rho$-RMSE spec & \texttt{group1/deep\_combined\_attn/L4\_balanced\_vxc\_anchor/full\_3} \\
random-init NN AE (kcal/mol) & 461.5 \\
pretrained-only NN AE (kcal/mol) & 121.4 \\
PBE AE-MAE on H$_2$O+C$_2$H$_2$ (kcal/mol) & 8.17 \\
random / best & $5.82\!\times\!10^{4}$ \\
pretrained-only / best & $1.53\!\times\!10^{4}$ \\
PBE / best & $1.03\!\times\!10^{3}$ \\
\bottomrule
\end{tabular}
\end{table}

\noindent
The unweighted-origin best spec lands at sub-millihartree AE
(0.00793 kcal/mol $\approx 1.26\!\times\!10^{-5}$ Ha) on its specific training
configuration. Per-loss group-3 MAEs: \texttt{L1\_B} 0.34, \texttt{L2\_C\_anchor}
0.34, \texttt{L3\_balanced\_vxc} 1.62, \texttt{L4\_balanced\_vxc\_anchor} 2.15,
\texttt{L5\_gradnorm\_vxc} 0.43 (kcal/mol). Anchor-augmented losses
(\texttt{L2\_C\_anchor}) match the unconstrained baseline (\texttt{L1\_B})
on AE; the explicit $V_{xc}$-aware losses
(\texttt{L3\_balanced\_vxc}, \texttt{L4\_balanced\_vxc\_anchor}) trade AE for
$V_{xc}$-fidelity.

\section{Post-processing figures}

Figures generated by
\texttt{notebooks/analysis/run\_step6\_unweighted\_analysis.py}.

\begin{figure}[h!]
\centering
\includegraphics[width=\textwidth]{fig1_baseline_reduction.png}
\caption{Baseline AE-error reduction ladder. Random ($\sim$461 kcal/mol),
pretrained-only NN (121 kcal/mol), PBE AE-MAE on the training set H$_2$O+C$_2$H$_2$
($\sim$8.17 kcal/mol; per-mol PBE errors are 7.03~kcal/mol on H$_2$O alone and
9.89~kcal/mol on C$_2$H$_2$ at this run's def2-svp / grid\_level=1 configuration,
\emph{not} a Karton 2011 literature value), chemical accuracy 1~kcal/mol
\citep{Pople1999}, and best NN spec.}
\end{figure}

\begin{figure}[h!]
\centering
\includegraphics[width=\textwidth]{fig1b_baseline_reduction_transfer.png}
\caption{Same ladder evaluated on the two transfer sets:
top row = primary $\{$CH$_4$, H$_2$, OH$\}$ (3 mols);
bottom row = secondary $\{$CO$_2$, HF, NH$_2$, NH$_3\}$ (4 mols);
plotting layout from \texttt{notebooks/analysis/comparison\_lib.py:360--453}.
PBE references are hardcoded for this specific def2-svp / grid\_level=1
configuration; \citet{KartonDaonMartin2011} W4-11 tabulates
CCSD(T)-quality reference atomization energies, not PBE errors, so direct
comparison to W4-11 requires running PBE on the W4-11 geometries.}
\end{figure}

\begin{figure}[h!]
\centering
\includegraphics[width=\textwidth]{fig2_density_vs_energy_tradeoff.png}
\caption{Density-vs-energy tradeoff. Each point is one spec; color = loss.
The L3/L4 cluster (explicit $V_{xc}$-balancing) sits to the right (worse AE)
but downward (tighter $\rho$-RMSE), exhibiting the classic Burke-Perdew
energy-vs-density tradeoff \citep[paraphrased in][]{Kepp2017,Burke1998}.}
\end{figure}

\begin{figure}[h!]
\centering
\includegraphics[width=\textwidth]{fig3_in_dist_vs_transfer.png}
\caption{In-distribution (training) AE vs.\ transfer-mol AE per spec. Loss
choice mostly does not regularize transfer: \texttt{L2\_C\_anchor} maintains
in-distribution AE without sacrificing transfer, but no spec reaches chemical
accuracy on the full W4-11 transfer set, mirroring the step-5 pattern
(\citet{BehlerParrinello2007}).}
\end{figure}

\begin{figure}[h!]
\centering
\includegraphics[width=\textwidth]{fig4_loss_strategy_heatmap.png}
\caption{Loss-strategy AE-MAE heatmap. Rows = loss family;
columns = the 18 (group, arch, solver) cells on training mols (H$_2$O+C$_2$H$_2$).
Cell value = mean $|\!$AE error$|\!$; colour = $\log_{10}$ of that mean
(\texttt{notebooks/analysis/comparison\_lib.py:636--674}).
Cells below chemical accuracy ($\log_{10} < 0$) appear only for
\texttt{L1\_B} / \texttt{L2\_C\_anchor} / \texttt{L5\_gradnorm\_vxc} on group-3
with the descriptor-rich \texttt{deep\_combined\_attn} arch.}
\end{figure}

\begin{figure}[h!]
\centering
\includegraphics[width=\textwidth]{fig5_arch_comparison.png}
\caption{Architecture comparison across the 5 loss families.
Bars (panel a, ``trained mols'') plot the \emph{mean} $|\!$AE error$|\!$ over all
9~(group, solver) cells per (loss, arch); see
\texttt{notebooks/analysis/comparison\_lib.py:681--752}.
Under that aggregation, \texttt{deep\_combined\_attn} wins on
\texttt{L1\_B}, \texttt{L2\_C\_anchor}, \texttt{L4\_balanced\_vxc\_anchor};
\texttt{deep\_combined} wins on \texttt{L3\_balanced\_vxc} and
\texttt{L5\_gradnorm\_vxc} (group-1/group-2 specs not yet converged with attention
inflate the attn means on those two losses).
Restricting to the converged group-3 cells flips the comparison to
\texttt{deep\_combined\_attn} winning all 5 by median MAE; per-loss group-3 medians
(deep vs.\ attn, kcal/mol):
\texttt{L1\_B} 0.65/0.04, \texttt{L2\_C\_anchor} 0.65/0.01,
\texttt{L3\_balanced\_vxc} 2.51/1.23,
\texttt{L4\_balanced\_vxc\_anchor} 2.58/1.53,
\texttt{L5\_gradnorm\_vxc} 0.83/0.02
(recomputed from \texttt{notebooks/checkpoints\_step6/unweighted/eval\_df.csv};
not the statistic the bars plot).}
\end{figure}

\section{In-notebook diagnostic figures}

The step-6 notebook (\texttt{notebooks/\_build\_step6\_notebook.py})
produces the following diagnostics during the run, independent of the
post-processing analysis. They cover (i) pretrain dynamics
(\texttt{pretrain\_losses.png}), (ii) anchor and balancing-loss
sensitivity studies (\texttt{anchor\_effect.png},
\texttt{balancing\_effect.png}), (iii) $F_x(s)$ profiles on
probe-molecule grids
(\texttt{fx\_drift\_panel\_B.png}, \texttt{fx\_drift\_panel\_C.png}),
(iv) per-group training-loss curves
(\texttt{loss\_curves\_group\{1,2,3\}\_*.png}), (v) SCF convergence
aggregate (\texttt{scf\_convergence.png}), (vi) transfer-set MAE
breakdowns (\texttt{transfer\_primary\_mae.png},
\texttt{transfer\_secondary\_mae.png}), and (vii) a $V_{xc}$-loss
efficacy diagnostic (\texttt{vxc\_efficacy.png}). Captions describe each
plot's semantic; the diagnostics are presented for completeness, not
as the headline results (see \S1--\S2 for those).

\begin{figure}[h!]
\centering
\includegraphics[width=0.92\textwidth]{pretrain_losses.png}
\caption{Pretrain loss curves per architecture (in-notebook).}
\end{figure}

\begin{figure}[h!]
\centering
\includegraphics[width=0.92\textwidth]{anchor_effect.png}
\caption{Anchor-loss effect on AE/density tradeoff (in-notebook).}
\end{figure}

\begin{figure}[h!]
\centering
\includegraphics[width=0.92\textwidth]{balancing_effect.png}
\caption{Balancing-loss effect (in-notebook).}
\end{figure}

\begin{figure}[h!]
\centering
\includegraphics[width=0.92\textwidth]{fx_drift_panel_B.png}
\caption{$F_x(s)$ sampled on the CH$_4$ (transfer-reference) and C$_2$H$_2$
(in-training) molecular grids; one subplot per probe molecule. Black =
analytic PBE eq.~14 \citep{PerdewBurkeErnzerhof1996}; green = per-arch
pretrained baselines; coloured scatter = all 90 trained models, one colour
per training group (G1/G2/G3). Plot generated by
\texttt{notebooks/\_build\_step6\_notebook.py:2935--3028}; not filtered by
loss family.}
\end{figure}

\begin{figure}[h!]
\centering
\includegraphics[width=0.92\textwidth]{fx_drift_panel_C.png}
\caption{$F_x(s)$ sampled on the C$_2$H$_4$ (ethylene) molecular grid --- a
held-out generalisation probe present in neither training nor transfer set.
Same colour key as the previous figure; all 90 trained models overlaid.
Plot generated by \texttt{notebooks/\_build\_step6\_notebook.py:3033--3134}.}
\end{figure}

\begin{figure}[h!]
\centering
\includegraphics[width=0.92\textwidth]{loss_curves_group1_h2o_short.png}
\caption{Loss curves for group1 (H$_2$O, short schedule).}
\end{figure}

\begin{figure}[h!]
\centering
\includegraphics[width=0.92\textwidth]{loss_curves_group2_h2o_c2h2_short.png}
\caption{Loss curves for group2 (H$_2$O + C$_2$H$_2$, short schedule).}
\end{figure}

\begin{figure}[h!]
\centering
\includegraphics[width=0.92\textwidth]{loss_curves_group3_h2o_c2h2_long.png}
\caption{Loss curves for group3 (H$_2$O + C$_2$H$_2$, long schedule).}
\end{figure}

\begin{figure}[h!]
\centering
\includegraphics[width=0.6\textwidth]{scf_convergence.png}
\caption{SCF-mode convergence aggregate (in-notebook).}
\end{figure}

\begin{figure}[h!]
\centering
\includegraphics[width=0.92\textwidth]{transfer_primary_mae.png}
\caption{Primary transfer-set MAE per spec (in-notebook).}
\end{figure}

\begin{figure}[h!]
\centering
\includegraphics[width=0.92\textwidth]{transfer_secondary_mae.png}
\caption{Secondary transfer-set MAE per spec (in-notebook).}
\end{figure}

\begin{figure}[h!]
\centering
\includegraphics[width=0.92\textwidth]{vxc_efficacy.png}
\caption{$V_{xc}$ efficacy diagnostic (in-notebook): correlation between
$V_{xc}$-loss-component magnitude and downstream density quality.}
\end{figure}

\section{Conclusions and recommendations (unweighted origin)}

\begin{enumerate}
\item \texttt{L1\_B} (in-distribution-only) and \texttt{L2\_C\_anchor}
  (in-distribution + close-shell anchor) co-lead on group-3 AE; they are
  statistically tied. Anchor adds no harm and improves $V_{xc}$-fidelity
  diagnostics.
\item \texttt{L3\_balanced\_vxc} and \texttt{L4\_balanced\_vxc\_anchor}
  produce the tightest density (best $\rho$-RMSE) at the cost of $\sim 5\times$
  worse AE than \texttt{L1\_B}/\texttt{L2\_C\_anchor}. This is the
  energy-vs-density tradeoff in action \citep[paraphrased in][]{Kepp2017,Burke1998};
  it is not a tuning failure --- it is the exact-functional approach showing up.
\item \texttt{L5\_gradnorm\_vxc} \citep{ChenGradNorm2018} sits in between:
  AE comparable to \texttt{L1\_B}/\texttt{L2\_C\_anchor}, $\rho$-RMSE
  comparable to \texttt{L3}/\texttt{L4}. The published GradNorm
  $\alpha\!=\!1.5$ was tuned for NYUv2 multitask vision and is not necessarily
  optimal for our $V_{xc}$-vs-AE balance --- a small grid sweep of $\alpha$
  is warranted.
\item Transfer to W4-11 \citep{KartonDaonMartin2011} is uniformly
  insufficient at this training-set size. Either expand training mol set
  (consistent with \citet{BehlerParrinello2007} for NN PESs) or use a
  larger/finer grid for $V_{xc}$ supervision.
\end{enumerate}

\bibliographystyle{plainnat}
\bibliography{references}
\end{document}
